\section{Step 31 Plan one request and six records}\label{h263-macroblock-s01}
Sketch the request offsets and the six output ranges before implementing the
routine. The request has eighteen doublewords. The output has six records of
192 bytes. A block's record starts at $192b$ for block index $b$; its signed
residuals start another 64 bytes later.

Keep saved caller state, a stable request copy, the reusable child request and
staged results distinct. Align the ordinary entry stack before its first call;
a private record at an aligned offset still needs an aligned stack base.
Record sizes are byte counts, while loop counters may
count doublewords. Confusing those units can produce a copy that looks correct
at its beginning and leaves the final Cr residuals unchanged.

\paragraph{Observe} Block five starts at byte 960 and its final halfword occupies
bytes 1150 and 1151. The next byte belongs to the output guard.
\Needspace{8\baselineskip}
\paragraph{Check} Where does block five's residual array start?
\begin{itemize}
  \item[A.] Byte 64, shared by every block.
  \item[B.] Byte 1024, which is $5\times192+64$.
  \item[C.] Byte 1152, after the whole result.
\end{itemize}

\section{Step 32 Check the macroblock geometry}\label{h263-macroblock-s02}
Check request/output alignment and extents, then load the coordinates,
displacement, mode and reserved field. Coordinates are unsigned and aligned to
sixteen samples. Displacements are signed full-width quantities. Validate their
small range before forming reference coordinates.

A valid current position does not make every displacement valid. At $x=0$,
$d_x=-1$ points outside the reference picture and must be rejected. A preceding
search would have omitted that candidate; this routine must still defend its
own interface when called directly.

\paragraph{Observe} Current origin $(16,32)$ is admitted. $(17,32)$ is not a
macroblock origin. $(0,32)$ with horizontal displacement $-1$ is rejected.
\Needspace{8\baselineskip}
\paragraph{Check} What should this routine do with an out-of-picture selected vector?
\begin{itemize}
  \item[A.] Reject it before pixel reads or output changes.
  \item[B.] Silently clamp it to a different vector.
  \item[C.] Repeat the complete motion search without being asked.
\end{itemize}

\section{Step 33 Validate every plane before starting Y}\label{h263-macroblock-s03}
Check all six plane strides and non-wrapping extents. Select width 176 and height
144 for Y, width 88 and height 72 for each chroma plane. Validate each current and
reference stride independently. A shared numerical value in the worked caller
is an example, not permission to ignore a field.

Finish this preflight before calling the first block predictor. Otherwise an
invalid reference Cr stride could be discovered only after Y has already been
read and written. Status one requires no pixel reads, even when the invalid
field occurs at the end of the request.

\paragraph{Observe} Setting only the final reference Cr stride to zero returns
one and zero, retaining the entire old output and leaving all image traces empty.
\Needspace{8\baselineskip}
\paragraph{Check} Which extent belongs to an 88-by-72 chroma plane?
\begin{itemize}
  \item[A.] $143s+176$, copied from Y.
  \item[B.] $72s+88$, with an extra row advance.
  \item[C.] $71s+88$, with 71 advances before the last active row.
\end{itemize}

\section{Step 34 Build the four luma child requests}\label{h263-macroblock-s04}
For blocks zero through three, derive local offsets $(0,0)$, $(8,0)$, $(0,8)$
and $(8,8)$. Add each to the current macroblock origin. The reference coordinate
also includes the selected displacement. Compute byte offsets with the respective
Y strides and complete plane origins.

Set both child phases to zero because this search selects full luma pixels.
Set the child's mode and reserved field explicitly. Retain all values needed for
later blocks across the call; child argument and scratch registers are temporary.
The outer return address must survive every nested call.

\paragraph{Observe} For current $(16,32)$ and displacement $(-2,-2)$, block one
uses current $(24,32)$ and reference $(22,30)$, both in Y.
\Needspace{8\baselineskip}
\paragraph{Check} Which local offset belongs to the bottom-left Y block?
\begin{itemize}
  \item[A.] $(8,0)$.
  \item[B.] $(0,8)$.
  \item[C.] $(8,8)$.
\end{itemize}

\section{Step 35 Derive both chroma requests from the same vector}\label{h263-macroblock-s05}
Halve the aligned current coordinates to obtain the chroma macroblock origin.
Split each signed luma displacement into its floored half and its phase bit.
Add the floored shifts to the chroma origin, then form addresses from the chosen
Cb or Cr plane and its own stride.

Reuse the geometry for both chroma blocks, but load their distinct pointers and
strides. An odd negative component needs origin shift $-1$ and phase one.
Truncating division toward zero would select the wrong reference interval.
The displayed sample aspect ratio does not change these array coordinates.

\paragraph{Observe} At current $(16,32)$ and displacement $(-1,-1)$, both chroma
blocks use current $(8,16)$, reference origin $(7,15)$ and phases $(1,1)$.
\Needspace{8\baselineskip}
\paragraph{Check} What may Cb and Cr legitimately share here?
\begin{itemize}
  \item[A.] The derived displacement geometry, while retaining their own data and strides.
  \item[B.] One pointer, regardless of the request.
  \item[C.] The Y stride, regardless of either chroma stride field.
\end{itemize}

\section{Step 36 Point the child at its private record}\label{h263-macroblock-s06}
Build the reusable eight-by-eight request from the current block's fields.
Pass a private output address at staged-base plus $192b$. The child writes its
prediction and residual there. Its successful return does not publish the whole
macroblock yet; the remaining records still have to be produced.

A stable block index identifies output order even when you process blocks in a
different sequence. An implementation that visits five through zero can store
each result at the same specified record offset. Do not make traversal order
implicitly determine which plane appears in a record.

\paragraph{Observe} After one successful child call, only its private record
is complete. The caller's output still has all of its original bytes.
\Needspace{8\baselineskip}
\paragraph{Check} Which address should the child receive as its output?
\begin{itemize}
  \item[A.] The caller output, before knowing whether later children succeed.
  \item[B.] The request's reference plane pointer.
  \item[C.] Its assigned private staging record.
\end{itemize}

\section{Step 37 Add magnitudes without cancelling errors}\label{h263-macroblock-s07}
On each successful child return, add its SAD to a full-width total. This total
summarizes all six records. Keep it separate from the byte-copy counter, block
index and status register. Save the returned status and total before an
ordinary observation caller prints them. Loading a syscall service number into
\texttt{\$v0} replaces the function's returned status.

Sum magnitudes, rather than taking the magnitude of a sum of signed residuals.
Residuals $-10$ and $+10$ contribute twenty to SAD even though their signed sum
is zero. The six-block upper bound is 97\,920, covering 384 current samples.
Every signed residual remains stored for the later transform.

\paragraph{Observe} With all current samples 255 and all reference samples zero,
the return total is 97\,920 in both modes. Reversing the extremes retains that
total but reverses the residual signs.
\Needspace{8\baselineskip}
\paragraph{Check} What do residuals $-10$ and $+10$ contribute to the returned total?
\begin{itemize}
  \item[A.] Zero, because they cancel.
  \item[B.] Twenty, the sum of their magnitudes.
  \item[C.] Two, the number of samples only.
\end{itemize}

\section{Step 38 Handle a failed child without publishing}\label{h263-macroblock-s08}
Inspect the child status before adding its score or proceeding to another block.
If it is nonzero, return dependency status two and total zero. Restore saved state
and leave the caller's output untouched. The child may have changed its private
record before reporting failure; those bytes are not a complete macroblock result.

A completed outer routine with the supplied unfinished child is a useful
integration observation. It should report block prediction failure, rather than
pretend the child's development value is a valid SAD. Restore the completed child
before expecting the ordinary success checks to pass.

\paragraph{Observe} A failure on a later child retains the old caller output,
even after earlier private records were successfully produced. Argument rejection
still uses status one and occurs before any child pixel work.
\Needspace{8\baselineskip}
\paragraph{Check} What does dependency status two imply about earlier input reads?
\begin{itemize}
  \item[A.] Earlier successful child calls may already have read their inputs.
  \item[B.] The function never read its request.
  \item[C.] The caller's first three output records must be overwritten.
\end{itemize}

\section{Step 39 Copy the complete staged output}\label{h263-macroblock-s09}
After all six children succeed, copy exactly 1152 bytes to the caller's output.
For aligned doubleword transfers this is 144 iterations. Keep a distinct source
and destination pointer and stop at the exclusive end. Preserve record order and
the native-endian halfword bytes already produced by the child.

Return success and the accumulated total only after the copy completes. The
caller must use the returned status to decide whether the output is new. This
working representation will feed transforms; it is not a stream of already
compressed H.263 bytes.

\paragraph{Observe} The default worked caller prints 2976. At its return
breakpoint, the sixth record is complete and the byte immediately after the
1152-byte area is unchanged.
\Needspace{8\baselineskip}
\paragraph{Check} How many eight-byte transfers cover the complete output?
\begin{itemize}
  \item[A.] Eighteen, the number of request fields.
  \item[B.] 1152, treating bytes as doublewords.
  \item[C.] 144.
\end{itemize}

\section{Step 40 Verify complete records before adding transforms}\label{h263-macroblock-s10}
Run completed feedback against your current source:
\begin{verbatim}
python3 construction/h263-macroblock-construction/grade.py \
    my-macroblock.asm --cli /path/to/mars-cli.jar \
    --report macroblock-results.json
\end{verbatim}
The 64 public cases cover both modes and byte orders, five signed vectors,
picture-corner macroblocks, padded independent strides, read-only aliases, maximum
differences and rejected requests. The checker computes every expected prediction
and signed residual from its own plane data. It checks all six record offsets,
the total, sample accesses, saved state, source/request preservation and output
guards. It accepts alternative block orders and private frames. Reports require
a new path.

Measure the outer entry-to-return boundary, including preflight, child setup,
all six calls, private staging and publication. For the named mixed-pattern
interior case with displacement $(-2,-2)$, the reference used 11\,447 scalar
instructions and 5525 with MSA block arithmetic. The sample results agreed.
These finite instruction counts do not measure host time or a completed encoder.
Publication and control work remain visible after faster block arithmetic.

\paragraph{Observe} Use the odd worked displacement as well as the default.
A correct total for a zero-chroma case cannot establish that Cr addressing and
half-sample chroma prediction are correct. Check their individual records.
\Needspace{8\baselineskip}
\paragraph{Check} Which boundary supplies the appropriate macroblock work count?
\begin{itemize}
  \item[A.] Only the arithmetic inside one child.
  \item[B.] Outer entry through return, including staging and publication.
  \item[C.] The number of bytes in the raw input file.
\end{itemize}
